\chapter{Ancient \texorpdfstring{$\kappa$}{kappa}-solutions}
\label{chap:cz-ancient-kappa-solutions}
\section{Preliminaries}
\begin{proposition}[lower radius bound for necks at infinity]
  \label{prop:cz-neck-radius-lower-bound-at-infinity}
  \source{cao-zhu}
  \dcref{Proposition 6.1.1}
  \group{Ancient-Solution Preliminaries}

There exists a positive constant $\varepsilon_0=\varepsilon_0(n)$
such that every complete noncompact $n$-dimensional Riemannian
manifold $(M, g_{ij})$ of nonnegative sectional curvature has a
positive constant $r_0$ such that any $\varepsilon$-neck of radius
$r$ on $(M, g_{ij})$ with $\varepsilon \leq \varepsilon_0$ must
have $r\geq r_0$.
\end{proposition}
\begin{proof}
\sketch A contradictory sequence of shrinking necks at infinity yields
minimizing rays through all necks.  Toponogov comparison of the induced
Busemann level sets contradicts their radii tending to zero \cite{CZ05F}.
\end{proof}
\begin{proposition}[line splitting in pointed limits at infinity]
  \label{prop:cz-pointed-limit-at-infinity-splits-line}
  \source{cao-zhu}
  \dcref{Proposition 6.1.2}
  \group{Ancient-Solution Preliminaries}

Suppose $(M,g_{ij})$ is a complete $n$-dimensional Riemannian
manifold with nonnegative sectional curvature. Let $P\in M$ be
fixed, and $P_k\in M$ a sequence of points and $\lambda_k$ a
sequence of positive numbers with $d(P,P_k)\rightarrow +\infty$
and $\lambda_kd(P,P_k)\rightarrow +\infty$. Suppose also that the
marked manifolds $(M,\lambda_k^2g_{ij},P_k)$ converge in the
$C^{\infty}_{loc}$ topology to a Riemannian manifold
$\widetilde{M}$. Then the limit $\widetilde{M}$ splits
isometrically as the metric product of the form $\mathbb{R}\times
N$, where $N$ is a Riemannian manifold with nonnegative sectional
curvature.
\end{proposition}
\begin{proof}
\sketch Rescaled geodesics from $P$ through a sparse subsequence of the
$P_k$ converge to opposite rays forming a line.  Toponogov splitting
then gives $\widetilde M=\mathbb R\times N$ \cite{CZ05F, Mi, BGP}.
\end{proof}
\begin{lemma} [point selection at infinite asymptotic curvature ratio]
  \label{lem:cz-infinite-ascr-point-selection}
  \source{cao-zhu}
  \dcref{Lemma 6.1.3}
  \group{Ancient-Solution Preliminaries}

 Given a complete noncompact Riemannian manifold with
bounded curvature and with asymptotic scalar curvature ratio
$$
A=\limsup_{s\rightarrow+\infty}Rs^2=+\infty,
$$
we can find a sequence of points $x_j$ divergent to infinity, a
sequence of radii $r_j$ and a sequence of positive numbers
$\delta_j\rightarrow0$ such that
\begin{itemize}
\item[(i)] $R(x)\leq(1+\delta_j)R(x_j)$ for all $x$ in the ball
$B(x_j,r_j)$ of radius $r_j$ around $x_j$, \item[(ii)]
$r^2_jR(x_j)\rightarrow+\infty$, \item[(iii)]
$\lambda_j=d(x_j,O)/r_j\rightarrow+\infty$, \item[(iv)]
the balls $B(x_j,r_j)$ are disjoint, \\
where $d(x_j,O)$ is the distance of $x_j$ from the origin $O$.
\end{itemize}
\end{lemma}
\begin{proof}
\sketch Select growing maxima of $Rs^2$, choose almost-maximal
curvature points outside them, and take radii proportional to their
distances from the origin.  The defining inequalities give (i)--(iii);
a subsequence gives (iv) \cite{Ha95F}.
\end{proof}
\begin{lemma}[point selection for unbounded curvature]
  \label{lem:cz-unbounded-curvature-point-selection}
  \source{cao-zhu}
  \dcref{Lemma 6.1.4}
  \group{Ancient-Solution Preliminaries}

Given a complete noncompact Riemannian manifold with unbounded
curvature, we can find a sequence of points $x_j$ divergent to
infinity such that for each positive integer $j$, we have
$|Rm(x_{j})|\geq j$, and
$$
|Rm(x)|\leq4|Rm(x_{j})|
$$
for $x\in B(x_j,\frac{j}{\sqrt{|Rm(x_{j})|}})$.
\end{lemma}
\begin{proof}
\sketch Repeatedly replace a point by one of at least four times its
curvature within the indicated radius.  The curvature growth and bounded
total displacement force termination, and the final point has the claim.
\end{proof}
\section{Asymptotic Shrinking Solitons}
\begin{theorem}[asymptotic shrinking soliton of an ancient kappa-solution]
  \label{thm:cz-ancient-kappa-asymptotic-shrinker}
  \source{cao-zhu}
  \dcref{Theorem 6.2.1}
  \group{Asymptotic Shrinking Solitons}

Let $g_{ij}(\cdot,t), -\infty<t<T$ with some $T>0$, be a nonflat
ancient $\kappa$-solution for some $\kappa >0$. Then there exist a
sequence of points $q_k$ and a sequence of times $t_k \rightarrow
-\infty$ such that the scalings of $g_{ij}(\cdot,t)$ around $q_k$
with factor $|t_k|^{-1}$ and with the times $t_k$ shifting to the
new time zero converge to a nonflat gradient shrinking soliton in
$C_{\rm loc}^\infty$ topology.
\end{theorem}
\begin{proof}
  \uses{lem:cz-reduced-distance-curvature-parabolic-annulus}
\sketch Choose reduced-distance minimizers at $\tau_k\to\infty$.
The annular estimates and noncollapsing give a smooth limit; reduced-volume
monotonicity makes its soliton defect vanish.  A flat limit contradicts
the reduced-volume bound \cite{P1}.
\end{proof}
\begin{theorem} [classification of ancient two-dimensional kappa-solutions]
  \label{thm:cz-two-dimensional-ancient-kappa-classification}
  \source{cao-zhu}
  \dcref{Theorem 6.2.2}
  \group{Asymptotic Shrinking Solitons}

The only nonflat ancient $\kappa$-solutions to Ricci flow on
two-dimensional manifolds are the round sphere $\mathbb{S}^2$ and
the round real projective plane $\mathbb{R}\mathbb{P}^2$.
\end{theorem}
\begin{proof}
  \uses{thm:cz-ancient-kappa-asymptotic-shrinker, lem:cz-unbounded-curvature-point-selection, prop:cz-pointed-limit-at-infinity-splits-line, prop:cz-compact-shrinking-surface-soliton-rigidity, thm:cz-positive-euler-surface-convergence}
\sketch The asymptotic shrinker has bounded curvature and is compact;
otherwise point selection and splitting give a forbidden nonflat limit.
It is therefore round, and the compact-flow convergence result identifies
the original solution \cite{Ha95F, P2}.
\end{proof}
\section{Curvature Estimates via Volume Growth}
\begin{lemma} [vanishing asymptotic volume ratio of ancient solutions]
  \label{lem:cz-ancient-solution-asymptotic-volume-ratio-zero}
  \source{cao-zhu}
  \dcref{Lemma 6.3.1}
  \group{Curvature Estimates from Volume Growth}

Let $M$ be an $n$-dimensional complete noncompact Riemannian
manifold. Suppose $g_{ij}(x,t)$, $x\in M$ and $t\in(-\infty,T)$
with $T>0$, is a nonflat ancient solution of the Ricci flow with
nonnegative curvature operator and bounded curvature. Then the
asymptotic volume ratio of the solution metric satisfies
$$
\nu_M(t)=\lim_{r\rightarrow+\infty}\frac{\Vol_t(B_t(O,r))}{r^n}=0
$$
for each $t\in(-\infty,T)$.
\end{lemma}
\begin{proof}
  \uses{thm:cz-two-dimensional-ancient-kappa-classification, lem:cz-infinite-ascr-point-selection, prop:cz-pointed-limit-at-infinity-splits-line}
\sketch Induct on dimension.  Positive asymptotic volume ratio and
point selection yield either a lower-dimensional nonflat ancient factor
or a nonflat local metric-cone limit, both contradictions.  The zero
asymptotic scalar-curvature ratio case similarly forces flatness
\cite{P1}.
\end{proof}
\begin{lemma}[parabolic point-picking curvature control]
  \label{lem:cz-parabolic-point-picking-curvature-control}
  \source{cao-zhu}
  \dcref{Lemma 6.3.2}
  \group{Curvature Estimates from Volume Growth}

For any positive constants $B, C$ with $B>4$ and $C>1000$, there
exists $1\leq A<\min\{\frac{1}{4}B,\frac{1}{1000}C\}$ which tends
to infinity as $B$ and $C$ tend to infinity and satisfies the
following property. Suppose we have a (not necessarily complete)
solution $g_{ij}(t)$ to the Ricci flow, defined on $M\times
[-t_0,0]$, so that at each time $t\in [-t_0,0]$ the metric ball
$B_t(x_0,1)$ is compactly contained in $M$. Suppose there exists a
point $(x',t') \in M\times (-t_0,0]$ such that
$$
d_{t'}(x',x_0)\leq \frac{1}{4} \; \mbox{ and }\;
|Rm(x',t')|>C+B(t'+t_0)^{-1}.
$$
Then we can find a point $(\bar{x},\bar{t}) \in M\times (-t_0,0]$
such that
$$
d_{\bar{t}}(\bar{x},x_0)<\frac{1}{3}\; \mbox{ with }\; Q
=|Rm(\bar{x},\bar{t})|>C+ B(\bar{t}+t_0)^{-1},
$$
and
$$
|Rm(x,t)|\leq 4Q
$$
for all $(-t_0<) \ \bar{t}-AQ^{-1}\leq t\leq \bar{t}$ and
$d_t(x,\bar{x})\leq \frac{1}{10}A^{\frac{1}{2}}Q^{-\frac{1}{2}}.$
\end{lemma}
\begin{proof}
  \uses{lem:cz-finite-parabolic-curvature-point-selection}
\sketch Apply the finite point-selection lemma, then use distance
distortion under its wider curvature bound to obtain the stated smaller
parabolic ball \cite{P1}.
\end{proof}
\begin{theorem}[curvature estimate from local volume growth]
  \label{thm:cz-local-volume-growth-curvature-estimate}
  \source{cao-zhu}
  \dcref{Theorem 6.3.3}
  \group{Curvature Estimates from Volume Growth}

For every $w>0$ there exist $B=B(w)<+\infty,\ \ C=C(w)<+\infty, \
\ \tau_0=\tau_0(w)>0,$ and $\xi=\xi(w)>0$ $($depending also on the
dimension$)$ with the following properties. Suppose we have a
$($not necessarily complete$)$ solution $g_{ij}(t)$ to the Ricci
flow, defined on $M\times[-t_0r_0^2,0],$ so that at each time
$t\in [-t_0r_0^2,0]$ the metric ball $B_t(x_0,r_0)$ is compactly
contained in $M.$
\begin{itemize}
\item[(i)] If at each time $t \in [-t_0r_0^2,0]$,
$$
Rm(.,t)\geq -r_0^{-2} \; \mbox{ on }\; B_t(x_0,r_0)
$$
$$
\mbox{and } \; \Vol_t(B_t(x_0,r_0))\geq wr_0^n,
$$
then we have the estimate
$$
|Rm(x,t)|\leq Cr_0^{-2}+B(t+t_0r_0^2)^{-1}
$$
whenever $-t_0r_0^2 < t\leq 0$ and $d_t(x,x_0)\leq
\frac{1}{4}r_0.$ \item[(ii)] If for some $0 < \bar{\tau} \leq
t_0$,
$$
Rm(x,t)\geq -r_0^{-2} \; \mbox{ for }\; t\in [-\bar{\tau}
r_0^2,0], x\in B_t(x_0,r_0),
$$
$$
\mbox{and } \; \Vol_0(B_0(x_0,r_0))\geq wr_0^n,
$$
then we have the estimates
$$
\Vol_t(B_t(x_0,r_0))\geq \xi r_0^n\; \mbox{ for all }\;
\max\{-\bar{\tau}r_0^2,-\tau_0r_0^2\}   \leq t \leq 0,
$$
and
$$
|Rm(x,t)|\leq Cr_0^{-2}+B(t -\max\{-\bar{\tau}r_0^2,-\tau_0r_0^2\}
)^{-1}
$$
whenever $\max\{-\bar{\tau}r_0^2,-\tau_0r_0^2\} < t\leq 0$ and
$d_t(x,x_0)\leq \frac{1}{4}r_0.$
\end{itemize}
\end{theorem}
\begin{proof}
  \uses{lem:cz-parabolic-point-picking-curvature-control, lem:cz-ancient-solution-asymptotic-volume-ratio-zero}
\sketch Scale to $r_0=1$.  A counterexample gives a point-picked
high-curvature region with a relative volume lower bound.  Rescaling and
compactness then produce a nonflat ancient solution of positive
asymptotic volume ratio, contradicting the preceding lemma.  Distance
distortion proves (ii) from (i) \cite{P1}.
\end{proof}
\section{Ancient \texorpdfstring{$\kappa$}{kappa}-solutions on Three-manifolds}
\begin{lemma}[classification of three-dimensional shrinking solitons]
  \label{lem:cz-three-dimensional-shrinking-soliton-classification}
  \source{cao-zhu}
  \dcref{Lemma 6.4.1}
  \group{Three-Dimensional Ancient Kappa-Solutions}
\index{classification of three-dimensional shrinking
solitons} Let $(M,g_{ij}(t))$ be a nonflat gradient shrinking
soliton on a three-manifold. Suppose $(M,g_{ij}(t))$ has bounded
and nonnegative sectional curvature and is $\kappa$-noncollapsed
on all scales for some $\kappa>0$. Then $(M,g_{ij}(t))$ is one of
the following:
\begin{itemize}
\item[(i)] the round three-sphere $\mathbb{S}^3$, or a metric
quotient of $\mathbb{S}^3$; \item[(ii)] the round infinite
cylinder $\mathbb{S}^2\times \mathbb{R}$, or one of its
$\mathbb{Z}_2$ quotients.
\end{itemize}
\end{lemma}
\begin{proof}
  \uses{thm:cz-three-manifold-positive-ricci-sphere, rem:cz-differentiable-sphere-subsequential-convergence, prop:cz-pointed-limit-at-infinity-splits-line, thm:cz-two-dimensional-ancient-kappa-classification}
\sketch A null direction splits the universal cover and yields a
cylindrical quotient.  Compact positive curvature is round.  The
noncompact positive-curvature case is ruled out by second variation,
splitting at infinity, and the Gauss--Bonnet level-set argument
\cite{P2}.
\end{proof}
\begin{proposition}[universal noncollapsing of ancient three-dimensional solutions]
  \label{prop:cz-ancient-three-dimensional-universal-noncollapsing}
  \source{cao-zhu}
  \dcref{Proposition 6.4.2}
  \group{Three-Dimensional Ancient Kappa-Solutions}
 There exists a positive constant $\kappa_0$ with
the following property. Suppose we have a nonflat
three-dimensional ancient $\kappa$-solution for some $\kappa >0$.
Then either the solution is $\kappa_0$-noncollapsed on all scales,
or it is a metric quotient of the round three-sphere.
\end{proposition}
\begin{proof}
  \uses{lem:cz-three-dimensional-shrinking-soliton-classification, thm:cz-ancient-kappa-asymptotic-shrinker, lem:cz-unbounded-curvature-point-selection, prop:cz-pointed-limit-at-infinity-splits-line, prop:cz-neck-radius-lower-bound-at-infinity, prop:cz-three-dimensional-positive-ricci-asymptotic-roundness}
\sketch Classify the asymptotic shrinker.  The spherical alternative
makes the original flow round; in the cylindrical case, the
reduced-volume argument supplies a universal local volume lower bound,
hence $\kappa_0$ \cite{P2}.
\end{proof}
\begin{theorem}[elliptic estimates for ancient kappa-solutions]
  \label{thm:cz-ancient-kappa-elliptic-estimates}
  \source{cao-zhu}
  \dcref{Theorem 6.4.3}
  \group{Three-Dimensional Ancient Kappa-Solutions}
 There exist a positive constant $\eta$ and a
positive increasing function $\omega:\quad [0,+\infty)\rightarrow
(0,+\infty)$ with the following properties. Suppose we have a
three-dimensional ancient $\kappa$-solution
$(M,g_{ij}(t)),-\infty<t\le 0,$ for some $\kappa>0$. Then
\begin{itemize}
\item[(i)]for every $x,y\in M$ and $t\in(-\infty,0]$, there holds
$$
R(x,t)\le R(y,t)\cdot \omega(R(y,t)d_t^2(x,y));
$$
\item[(ii)] for all $x\in M$ and $t\in (-\infty,0]$, there hold
$$
|\nabla R|(x,t)\le\eta R^\frac{3}{2}(x,t)\;  \mbox{ and }\;
\left|\frac{\partial R}{\partial t}\right|(x,t)\le \eta R^2(x,t).
$$
\end{itemize}
\end{theorem}
\begin{proof}
  \uses{thm:cz-local-volume-growth-curvature-estimate}
\sketch Universal noncollapsing, Harnack, and local volume-growth
curvature control give the comparison function $\omega$.  Shi estimates
then give the spatial and temporal derivative bounds \cite{P1}.
\end{proof}
\begin{corollary}[compactness of normalized ancient kappa-solutions]
  \label{cor:cz-ancient-kappa-solution-compactness}
  \source{cao-zhu}
  \dcref{Corollary 6.4.4}
  \group{Three-Dimensional Ancient Kappa-Solutions}
 The set of
nonflat three-dimensional ancient $\kappa_0$-solutions is compact
modulo scaling in the sense that for any sequence of such
solutions and marking points $(x_k,0)$ with $R(x_k,0)=1$, we can
extract a $C^\infty_{loc}$ converging subsequence whose limit is
also an ancient $\kappa_0$-solution.
\end{corollary}
\begin{proof}
  \uses{thm:cz-ancient-kappa-elliptic-estimates, lem:cz-unbounded-curvature-point-selection, prop:cz-pointed-limit-at-infinity-splits-line, thm:cz-two-dimensional-ancient-kappa-classification, lem:cz-three-dimensional-shrinking-soliton-classification, prop:cz-neck-radius-lower-bound-at-infinity}
\sketch Elliptic estimates and noncollapsing give a local smooth limit.
Unbounded limiting curvature would rescale to a cylindrical limit,
contradicting the lower neck-radius bound.
\end{proof}
\begin{corollary} [cap-or-neck structure of noncompact ancient solutions]
  \label{cor:cz-noncompact-ancient-cap-neck-structure}
  \source{cao-zhu}
  \dcref{Corollary 6.4.5}
  \group{Three-Dimensional Ancient Kappa-Solutions}

For any $\varepsilon>0$ there exists $C=C(\varepsilon)>0$, such
that if $g_{ij}(t)$ is a nonflat ancient $\kappa$-solution on a
noncompact three-manifold $M$ for some $\kappa>0$, and
$M_\varepsilon$ denotes the set of points in $M$ which are not
centers of evolving $\varepsilon$-necks at $t=0$, then at $t=0$,
either the whole manifold $M$ is the round cylinder $\mathbb{S}^2
\times \mathbb{R}$ or its $\mathbb{Z}_2$ metric quotients, or
$M_\varepsilon$ satisfies the following
\begin{itemize}
\item[(i)] $M_\varepsilon$ is compact, \item[(ii)] ${\rm diam}\,
M_\varepsilon \le CQ^{-\frac{1}{2}}$ and $C^{-1}Q\le R(x,0)\le
CQ$, whenever $x\in M_\varepsilon$, where $Q=R(x_0,0)$ for some
$x_0\in\partial M_\varepsilon$.
\end{itemize}
\end{corollary}
\begin{proof}
  \uses{lem:cz-three-dimensional-shrinking-soliton-classification, prop:cz-pointed-limit-at-infinity-splits-line, thm:cz-two-dimensional-ancient-kappa-classification, cor:cz-ancient-kappa-solution-compactness, prop:cz-ancient-three-dimensional-universal-noncollapsing, thm:cz-ancient-kappa-elliptic-estimates}
\sketch A null direction gives the cylindrical cases.  Otherwise a
sequence of non-neck centers escaping to infinity has a normalized
limit that is cylindrical, a contradiction.  Elliptic estimates give
the compact-core diameter and curvature control \cite{P1}.
\end{proof}
\begin{theorem}[canonical neighborhoods in ancient kappa-solutions]
  \label{thm:cz-ancient-kappa-canonical-neighborhoods}
  \source{cao-zhu}
  \dcref{Theorem 6.4.6}
  \group{Three-Dimensional Ancient Kappa-Solutions}
 For any $\varepsilon>0$ one can find
positive constants $C_1=C_1(\varepsilon)$ and
$C_2=C_2(\varepsilon)$ with the following property. Suppose we
have a three-dimensional nonflat $($compact or noncompact$)$
ancient $\kappa$-solution $(M,g_{ij}(x,t))$. Then either the
ancient solution is the round $\mathbb{R}\mathbb{P}^2 \times
\mathbb{R}$, or every point $(x,t)$ has an open neighborhood $B$,
with $B_t(x,r)\subset B \subset B_t(x,2r)$ for some
$0<r<C_1R(x,t)^{-\frac{1}{2}}$, which falls into one of the
following three categories:
\begin{itemize}
\item[(a)] $B$ is an {\bf evolving $\varepsilon$-neck} $($in the
sense that it is the slice at the time $t$ of the parabolic region
$\{(x',t')\ |\ x' \in B, t' \in [t - \varepsilon^{-2}R(x,t)^{-1},
t] \}$ which is, after scaling with factor $R(x,t)$ and shifting
the time $t$ to zero, $\varepsilon$-close $($in the
$C^{[\varepsilon^{-1}]}$ topology$)$ to the subset
$(\mathbb{S}^2\times \mathbb{I})\times [-\varepsilon^{-2},0]$ of
the evolving standard round cylinder with scalar curvature $1$ and
length $2\varepsilon^{-1}$ to $\mathbb{I}$ at the time zero$),$ or
\index{evolving $\varepsilon$-neck} \item[(b)] $B$ is an {\bf
evolving $\varepsilon$-cap} $($in the sense that it is the time
slice at the time $t$ of an evolving metric on $\mathbb{B}^3$ or
$\mathbb{R}\mathbb{P}^3\setminus\bar{\mathbb{B}^3}$ such that the
region outside some suitable compact subset of $\mathbb{B}^3$ or
$\mathbb{R}\mathbb{P}^3\setminus\bar{\mathbb{B}^3}$ is an evolving
$\varepsilon$-neck$),$ or \index{evolving $\varepsilon$-cap}
\item[(c)] $B$ is a compact manifold $($without boundary$)$ with
positive sectional curvature $($thus it is diffeomorphic to the
round three-sphere $\mathbb{S}^3$ or a metric quotient of
$\mathbb{S}^3);$
\end{itemize}
furthermore, the scalar curvature of the ancient $\kappa$-solution
on $B$ at time $t$ is between $C_2^{-1}R(x,t)$ and $C_2R(x,t)$,
and the volume of $B$ in case {\rm (a)} and case {\rm (b)}
satisfies
$$
(C_2R(x,t))^{-\frac{3}{2}} \leq \Vol_t(B) \leq \varepsilon r^3.
$$
\end{theorem}
\begin{proof}
  \uses{thm:cz-two-dimensional-ancient-kappa-classification, lem:cz-three-dimensional-shrinking-soliton-classification, cor:cz-noncompact-ancient-cap-neck-structure, thm:cz-ancient-kappa-elliptic-estimates, prop:cz-ancient-three-dimensional-universal-noncollapsing, cor:cz-ancient-kappa-solution-compactness}
\sketch The noncompact case follows from the cap-or-neck structure.
A failing compact normalized sequence has a noncompact smooth limit,
whose cap-or-neck structure pulls back to a contradiction.  The null
and round cases are direct \cite{P2, CZ05F}.
\end{proof}
\section{Named intermediate formalization obligations}
\begin{lemma}[uniform reduced-distance and curvature bounds on parabolic annuli]
  \label{lem:cz-reduced-distance-curvature-parabolic-annulus}
  \source{cao-zhu}
  \dcref{Claim (6.2.6)}
  \group{Asymptotic Shrinking Solitons}
We first claim that for any $A \geq 1$, one can find
$B=B(A)<+\infty$ such that for every $\bar{\tau}>1$ there holds
\begin{equation} l(q,\tau)\leq B \; \mbox{ and }\;  \tau R(q,t_0-\tau)\leq B,
%\eqno(6.2.6)
\end{equation} whenever $\frac{1}{2}\bar{\tau}\leq \tau \leq A\bar{\tau}$ and
$d_{t_0-\frac{\bar{\tau}}{2}}^2 (q,q(\frac{\bar{\tau}}{2}))\leq
A\bar{\tau}.$

\end{lemma}
\begin{proof}
\sketch The gradient estimate for reduced distance bounds $l$ at
$\bar\tau/2$; Harnack bounds curvature on the annulus, and the
reduced-distance evolution inequality propagates both bounds.
\end{proof}
\begin{lemma}[finite parabolic curvature point selection]
  \label{lem:cz-finite-parabolic-curvature-point-selection}
  \source{cao-zhu}
  \dcref{Assertion (6.3.4)}
  \group{Curvature Estimates from Volume Growth}
We first claim that there exists a point $(\bar{x},\bar{t})$ with
$-t_0<\bar{t}\leq 0$ and $d_{\bar{t}}(\bar{x},x_0)<\frac{1}{3}$
such that
$$
Q=|Rm(\bar{x},\bar{t})|>C+ B(\bar{t}+t_0)^{-1},
$$
and \begin{equation}
|Rm(x,t)|\leq 4Q %\eqno (6.3.1)? ---> 6.3.4     WRONG number
\end{equation} wherever $\bar{t}-AQ^{-1}\leq t\leq \bar{t},\ \ d_t(x,x_0)\leq
d_{\bar{t}}(\bar{x},x_0) +(AQ^{-1})^{\frac{1}{2}}.$
\end{lemma}
\begin{proof}
\sketch Repeatedly replace a failing point by one of four times the
curvature in the permitted parabolic region.  Geometric-series bounds
on displacement and backward time, together with smoothness, force a
finite terminal point.
\end{proof}
